\chapter*{如何使用這本書}
\addcontentsline{toc}{chapter}{如何使用這本書}

如果你從來沒寫過程式，這本書就是為你寫的。不管你現在是國中生、高中生、大學生，還是已經
在職場工作好幾年的社會人士，這本書都預設你完全沒有寫過程式，也沒有用過任何特定的辦公
軟體或開發工具。書裡需要打比方的地方，會用大家都有的生活經驗（食譜步驟、手機 App 怎麼
運作、日常生活裡的規則跟流程），不會假設你用過某個特定軟體。

這本書有兩個並重的目標，不是只有一個。第一個目標，是讓你能自己動手寫小程式，解決生活、
課業或工作上的瑣事，例如整理一堆檔案、自動化重複性的工作、做簡單的資料分析。終點不是
「精通 Python 語言的每一個角落」，而是看到一件瑣碎的任務時，會反射性想到「這個我可以
寫程式做」，而且真的做得出來。

第二個目標，是讓你理解「軟體」這個東西究竟是怎麼被打造出來的：版本控制在做什麼、一個
專案的檔案通常長什麼樣子、前端、後端、API、部署這些詞在講什麼、工程師團隊的日常運作
方式大概是什麼樣子。這個目標特別是為了讓沒有工程背景、但工作上需要跟工程師對接的人
（最典型的是 PM），有能力聽懂、問對問題，不會被術語唬住。第 \ref{ch:industry}、\ref{ch:api} 章是專門服務這個
目標的章節。

這本書不會把你訓練成職業工程師，也不是那種「系統性教完整個語言」或「給有經驗的工程師
精進」的書。如果你剛好正卡在某堂必修的程式設計課，這本書當然也能幫上忙，但這不是這本書
設定的主要對象，你不會在書裡看到「這在 C 語言裡怎麼寫」這種對照內容。

\section*{每一節大致的節奏}

大部分小節都會照著同一個節奏走：先用一段白話文字說明這個技巧是什麼、為什麼需要它，
接著給你一個通用的片語樣板，然後是一支真正能執行的完整範例程式，範例後面緊接著這支
程式實際的輸入輸出。有些地方會補上一段特別提醒（標成「學習要點」或「風格要點」），
挑的都是容易跟你原本的直覺互相衝突、或是比較道地的 Python 寫法。

很多小節還會示範一次新手常見的錯誤寫法，讓你看看那個錯誤訊息長什麼樣子、該怎麼讀懂它、
怎麼修正。打字打錯一個字元，程式就整個不動，這件事任何人第一次寫程式都會碰到，跟你是
國中生還是已經工作多年沒有關係。看到紅色的錯誤訊息不代表你不是這塊料，只代表 Python
想告訴你哪裡需要修改，讀懂它、修好它，你就往前進了一步。

大部分小節最後會有兩到四題簡短的練習，讓你動手確認自己真的懂了；每一章的最後則集中
放著全章練習的解答，方便你對答案。

\section*{書裡會用到的幾種方塊}

程式碼樣板放在深色標題的方塊裡，標題是這個樣板的用途；真正能執行的完整範例也用類似的
方塊呈現，標題會附上檔名。範例後面常常接著一個標著「input」和「output」的小方塊，那是
這支程式實際的輸入跟輸出內容。「學習要點」跟「風格要點」是灰底的提示段落，通常一句話
就講完一個重點，值得多看一眼。

\section*{兩個會反覆出現的思考方式}

這本書會用兩個貫穿全書的思考方式，幫你建立對 Python 的直覺，而不是要你死背一堆規則。

第一個是「名字牌模型」：把變數想成一張貼在物件上的名字牌，而不是一個裝著值的盒子。這個
想法第 \ref{ch:variables} 章會第一次完整說明，之後在講清單、函式、tuple 的章節都會重新呼應一次，因為
每次呼應的情境不太一樣，能看到的東西也不一樣。

第二個是「先做、錯了再處理」（EAFP），跟「先檢查、確定安全才動手」（LBYL）相對照。
這個想法第 \ref{ch:files-exceptions} 章講檔案與例外處理時會第一次完整說明，之後在講字典、除錯的章節也會
重新呼應。

第一次看到這兩個名字，不用急著完全弄懂，書裡每次呼應時都會簡短提醒一次背景，多看幾次
自然就有感覺了。

\section*{建議的讀法}

這本書大致依照由簡入繁的順序安排，前面幾章（尤其是第 \ref{ch:variables} 到 \ref{ch:files-exceptions} 章）是後面所有內容的
地基，建議照順序讀完，不要跳著看。第 \ref{ch:csv} 章之後的章節彼此之間的關聯性比較低，如果你
已經有明確想解決的問題（例如想學怎麼串接 API），可以考慮直接跳到對應的章節。

不用急著一次讀完一整章，每個小節後面的練習題值得你真的動手做一次，而不是只看過去就好。
程式是動手做出來的技能，不是背出來的知識。

祝你讀得順利，也祝你很快就能親手寫出第一支真正解決問題的程式。
